\chapter{A Concrete Naming Certificate for $\PointwiseContinuityModulusUp$}
\label{ch:concrete-instances-pointwise-continuity-modulus-namecert}
\origin{human}

Following Bishop--Bridges pointwise continuity on metric spaces, the $\PointwiseContinuityModulusUp$ packet records the local radius surface that a displayed continuous map supplies at a single named source point and target tolerance. It sits between the metric source of \autoref{ch:concrete-instances-metric-namecert}, the graph and pointwise continuity rows of \autoref{ch:concrete-instances-continuousmap-namecert}, the rational tolerance rows of \autoref{ch:concrete-instances-rat-namecert}, the global modulus interface of \autoref{ch:concrete-instances-uniform-continuity-modulus-namecert}, and the compact finite-net extraction of \autoref{ch:concrete-instances-compact-uniform-continuity-modulus-namecert}. The packet names one point-indexed local radius ledger; it is not a uniform modulus, not a compactness principle, not a selected Lebesgue number, and not host semantic continuity.

\begin{definition}[Pointwise continuity modulus carrier]
\label{def:pointwise-continuity-modulus-carrier}
A $\PointwiseContinuityModulusUp$ carrier is a finite $\BHist$ packet
$$
P=(X,Y,F,x,\epsilon,\delta,R,U,K,H,C,Q,N)
$$
with the following displayed rows:
\begin{enumerate}
\item $X$ is the source $\MetricUp$ row whose point classifier and distance comparisons supply the neighbourhood side of the local read.
\item $Y$ is the target $\MetricUp$ row whose distance comparisons supply the tolerance side.
\item $F$ is the $\ContinuousMapUp$ graph row from $X$ to $Y$, including the displayed pointwise continuity evidence read at the named point.
\item $x$ is the finite source-point name at which the local radius is certified.
\item $\epsilon$ is the requested $\RatUp$ target tolerance row.
\item $\delta$ is the offered $\RatUp$ source radius row.
\item $R$ is the local radius ledger tying $x,\epsilon,\delta$ to the pointwise continuity row carried by $F$.
\item $U$ is the optional $\UniformContinuityModulusUp$ comparison row, read only when a downstream consumer asks whether the local radius is compatible with a global modulus surface.
\item $K$ is the optional $\CompactUniformContinuityModulusUp$ extraction row, read only when compact finite-net folding consumes the point-indexed radius surface.
\item $H$ records componentwise $\hsame$ transport for the two metric rows, graph row, point row, tolerance row, radius row, comparison rows, replay, provenance, and local naming.
\item $C$ records displayed $\Cont$ replay from a target-tolerance request through the local radius ledger and onward to any global or compact modulus consumer.
\item $Q$ is the $\Pkg$ provenance over $\PointwiseContinuityModulusUp$, $\MetricUp$, $\ContinuousMapUp$, $\RatUp$, $\UniformContinuityModulusUp$, $\CompactUniformContinuityModulusUp$, $\BHist$, $\hsame$, $\Cont$, and $\NameCert$.
\item $N$ is the local $\NameCert_{\PointwiseContinuityModulusUp}$ row.
\end{enumerate}
The classifier compares accepted carriers only through $X,Y,F,x,\epsilon,\delta,R,U,K,H,C,Q,N$. It has no coordinate for a host neighbourhood system, a selected global modulus function, compact open-cover data, sequence extraction, countable choice, or ambient completeness.
\end{definition}

\begin{theorem}[Pointwise continuity modulus NameCert obligations]
\label{thm:pointwise-continuity-modulus-namecert-obligations}
For an accepted $\PointwiseContinuityModulusUp$ carrier $P=(X,Y,F,x,\epsilon,\delta,R,U,K,H,C,Q,N)$, the local $\NameCert_{\PointwiseContinuityModulusUp}$ surface consists exactly of carrier admission for the displayed rows, source-target metric lock through $X$ and $Y$, graph compatibility through $F$, point-indexed tolerance visibility through $x$ and $\epsilon$, radius exactness through $\delta$ and $R$, optional global-modulus compatibility through $U$, optional compact-extraction compatibility through $K$, structural transport through $H$, displayed replay through $C$, provenance through $Q$, and local naming through $N$.
\end{theorem}
\begin{proof}
Project the carrier of \autoref{def:pointwise-continuity-modulus-carrier}. The only source of a local radius is the displayed tuple $x,\epsilon,\delta,R$ attached to the continuous-map row $F$ over the two metric rows $X$ and $Y$.

The ledger $R$ binds the requested target tolerance to the offered source radius at the named point. Thus the packet records a local point-indexed read, not a uniform radius for all source points and not a host-level continuity assertion.

The optional rows $U$ and $K$ are downstream comparison surfaces. A $\UniformContinuityModulusUp$ consumer may read $U$ only after the local radius ledger has been displayed, and the compact extraction route of \autoref{ch:concrete-instances-compact-uniform-continuity-modulus-namecert} may read $K$ only as a compact finite-net consumer of the same point-indexed ledger.

Rows $H,C,Q,N$ transport, replay, source, and name the same finite route. Since no host neighbourhood system, selected global modulus function, compact open-cover datum, sequence extraction principle, countable-choice row, or ambient-completeness row occurs in the carrier, the listed obligations exhaust the NameCert surface.
\end{proof}

\begin{theorem}[Pointwise radius handoff to uniform extraction]
\label{thm:pointwise-continuity-modulus-uniform-extraction-handoff}
Every accepted $\PointwiseContinuityModulusUp$ carrier exports a local-radius read through the finite route
$$
\MetricUp,\ \ContinuousMapUp,\ \RatUp,\ \UniformContinuityModulusUp,\ \CompactUniformContinuityModulusUp,\ \BHist,\ \hsame,\ \Cont,\ \Pkg,\ \NameCert .
$$
A downstream uniform-modulus or compact-uniform-modulus consumer receives the point row $x$, the requested target tolerance $\epsilon$, the offered source radius $\delta$, and the ledger $R$ before it may consult either $U$ or $K$. The route exports no selected global modulus function, no open-cover compactness, no selected Lebesgue number, no sequence extraction principle, and no ambient complete metric space.
\end{theorem}
\begin{proof}
Start with the obligation surface of \autoref{thm:pointwise-continuity-modulus-namecert-obligations}. It fixes the two metric rows, the continuous-map graph row, the point row, the requested tolerance, the offered radius, and the local ledger as the only pointwise source of modulus information.

Now read a uniform consumer. The row $U$ is available only as a comparison against the already displayed pointwise ledger, so it cannot choose a radius independently of $x,\epsilon,\delta,R$.

Finally read a compact extraction consumer. The compact-uniform-continuity route of \autoref{ch:concrete-instances-compact-uniform-continuity-modulus-namecert} consumes pointwise radius rows center by center before folding them into a compact finite-net output. Therefore $K$ is downstream of the same local ledger, and the excluded host compactness and global selection principles have no coordinate in the handoff.
\end{proof}

\closureat{PointwiseContinuityModulusUp}{seedStr}

\begin{closurestatus}{\PointwiseContinuityModulusUp}
  \theoryclosure{\seedClosure}
  \scopeclosed{\autoref{def:pointwise-continuity-modulus-carrier}, \autoref{thm:pointwise-continuity-modulus-namecert-obligations}, and \autoref{thm:pointwise-continuity-modulus-uniform-extraction-handoff} seed the finite local-radius route over metric rows, ContinuousMap graph rows, Rat tolerance rows, UniformContinuityModulus comparison rows, CompactUniformContinuityModulus extraction rows, transport, replay, provenance, and local NameCert rows.}
  \formalstatus{\unformalizedV}
  \bridgestatus{none}
  \notclaimed{This chapter does not close a global uniform modulus, compact finite-net extraction, selected Lebesgue numbers, host continuity, Lean verification, or bridge checking.}
  \upgradepath{Obligation closure requires checked carrier admission, metric lock, point-indexed radius exactness, UniformContinuityModulus compatibility, CompactUniformContinuityModulus handoff, structural replay, and ledger non-escape for BEDC.Derived.PointwiseContinuityModulusUp.}
  \constructivestory{A $\PointwiseContinuityModulusUp$ packet is generated from source and target MetricUp rows, a ContinuousMapUp graph row, one named source point, a RatUp target tolerance, a RatUp source radius, a local radius ledger, optional UniformContinuityModulusUp and CompactUniformContinuityModulusUp comparison rows, componentwise hsame transport, Cont replay, provenance, and local naming.}
\end{closurestatus}
